% F081: precise mathematical teaching diagram.
\begin{tikzpicture}[x=1cm,y=1cm,>=stealth,font=\small]
\draw[rounded corners=3pt,gray!40] (0,0) rectangle (14.2,8.6);
\begin{scope}[yshift=1cm]
\node at (7.1,7.0) {易错一：都有水平切线，原函数零点却可以是2、1、0个};
\node[] at (2.6,6.3) {$x^2-1$};
\begin{scope}[shift={(2.6,3.6)},x=1.1cm,y=1.0cm]
\draw[->,gray!65] (-1.65,0)--(1.65,0) node[right] {$x$};
\draw[->,gray!65] (0,-1.6)--(0,2) node[left] {$y$};
\draw[blue!75!black,line width=1.15pt,,domain=-1.35:1.35,samples=160] plot (\x,{(\x)*(\x)+(-1)});
\draw[orange!85!black,line width=1.2pt] (-0.32,-1)--(0.32,-1);
\fill[orange!85!black] (0,-1) circle (1.8pt);
\fill[blue!75!black] (-1,0) circle (1.8pt);
\fill[blue!75!black] (1,0) circle (1.8pt);
\end{scope}
\node[] at (2.6,1.15) {2个零点};
\node[] at (7.1,6.3) {$x^2$};
\begin{scope}[shift={(7.1,3.6)},x=1.1cm,y=1.0cm]
\draw[->,gray!65] (-1.65,0)--(1.65,0) node[right] {$x$};
\draw[->,gray!65] (0,-1.6)--(0,2) node[left] {$y$};
\draw[blue!75!black,line width=1.15pt,,domain=-1.35:1.35,samples=160] plot (\x,{(\x)*(\x)+(0)});
\draw[orange!85!black,line width=1.2pt] (-0.32,0)--(0.32,0);
\fill[orange!85!black] (0,0) circle (1.8pt);
\fill[blue!75!black] (0,0) circle (1.8pt);
\end{scope}
\node[] at (7.1,1.15) {1个零点};
\node[] at (11.6,6.3) {$x^2+1$};
\begin{scope}[shift={(11.6,3.6)},x=1.1cm,y=1.0cm]
\draw[->,gray!65] (-1.65,0)--(1.65,0) node[right] {$x$};
\draw[->,gray!65] (0,-1.6)--(0,2) node[left] {$y$};
\draw[blue!75!black,line width=1.15pt,,domain=-1.35:1.35,samples=160] plot (\x,{(\x)*(\x)+(1)});
\draw[orange!85!black,line width=1.2pt] (-0.32,1)--(0.32,1);
\fill[orange!85!black] (0,1) circle (1.8pt);
\end{scope}
\node[] at (11.6,1.15) {0个零点};
\end{scope}
\node[text width=13cm,align=center] at (7.1,.9) {三个函数都满足 $f^{\prime}(0)=0$；存在导数零点不能反推存在两个函数零点。};
\end{tikzpicture}
